\documentclass[11pt]{article}
\usepackage[margin=1in]{geometry}
\usepackage[T1]{fontenc}
\usepackage[utf8]{inputenc}
\usepackage{amsmath, amssymb}
\usepackage{hyperref}
\title{systemTools.wl Module Documentation}
\author{MFGraphs maintainers}
\date{}
\begin{document}
\maketitle

\section{High-level explanation}

\texttt{systemTools.wl} is the structural equation-system kernel of MFGraphs. It takes a typed scenario and a typed symbolic unknown bundle and builds a typed \texttt{mfgSystem} object containing the equations, inequalities, rules, and typed subrecords that represent a stationary mean field game on the network.

Architecturally, this module is where scenario semantics become algebra. It collects boundary conditions, flow balance, complementarity constraints, switching inequalities, and Hamiltonian relations into a consistent system object. The design is modular: four builders produce typed records for boundary, flow, complementarity, and Hamiltonian data, and those records are then merged into the final system association.

\section{The mean field game system this module assembles}

Fix a directed multigraph $G = (V, E)$ obtained by augmenting the user-supplied network with auxiliary entry and exit vertices. On each directed edge $(a, b) \in E$ the module introduces a flow variable $j[a, b]$ and a per-vertex value variable $u[a, b]$ representing the optimal cost-to-go evaluated at vertex $a$ along the edge $(a, b)$. Three structural blocks act on these unknowns.

\paragraph{Boundary conditions.} For each entry vertex $a$ with mass $\mu_a$, the module emits the entry-current equation
\[
\sum_{b\,:\,(a,b)\in E} j[a, b] \;=\; \mu_a,
\]
and for each exit vertex $a$ with terminal cost $c_a$ the substitution rule $u[\cdot, a] \mapsto c_a$. These exit substitutions are propagated through the rest of the system at build time, so the symbolic solver never has to discover them.

\paragraph{Flow balance (Kirchhoff).} At every interior original vertex $i \in V$, the splitting and gathering balances impose
\[
j[\cdot,i] \;=\; \sum_{w} j[i,\,i,\,w]
\qquad\text{(gathering)},
\qquad
j[r,\,i,\,\cdot] \;=\; \sum_{w} j[r,\,i,\,w]
\qquad\text{(splitting)},
\]
where the auxiliary triples $\{(r, i, w)\}$ enumerate admissible transitions through $i$. These are the Kirchhoff laws that say "flow in equals flow out", lifted to the auxiliary topology so that switching costs can be charged on $(r, i, w)$ rather than on $i$ alone.

\paragraph{Hamiltonian residual (edge-level HJB).} For every undirected edge $\{a, b\}$ the module records
\[
u[a, b] - u[b, a] + j[a, b] - j[b, a] \;=\; \mathrm{nrhs}_{ab},
\qquad
\mathrm{nrhs}_{ab} \;=\;
\begin{cases}
0 & \text{if } \alpha = 1 \text{ (critical regime)}, \\[2pt]
m - \operatorname{sign}(m)\, m^{\alpha} & \text{otherwise},
\end{cases}
\]
which is the integrated form of the edge-level Hamilton--Jacobi--Bellman equation. The factor $\operatorname{sign}(m)$ makes the right-hand side orientation-aware so that one formula handles both travel directions on the undirected edge. The active critical-congestion solver path is exactly the case $\mathrm{nrhs}_{ab} = 0$.

\paragraph{Complementarity (LCP).} Two families of disjunctive constraints encode the equilibrium condition that flow only exists on edges that are part of an optimal path. On every undirected physical edge $\{a, b\}$,
\[
j[a, b] = 0 \;\;\vee\;\; j[b, a] = 0
\qquad\text{(alternative zero-flow)},
\]
and on every admissible transition triple $(r, i, w)$,
\[
j[r, i, w] \;=\; 0 \;\;\vee\;\; u[w, i] + \mathrm{sc}(r, i, w) - u[r, i] \;=\; 0,
\]
paired with the switching optimality inequality $u[w, i] + \mathrm{sc}(r, i, w) - u[r, i] \ge 0$. Together they form the linear complementarity problem (LCP) that the symbolic solvers in \texttt{solversTools.wl} reduce.

\section{Low-level explanation of functions and functionality}

\subsection*{\texttt{mfgSystem}}

\texttt{mfgSystem[assoc]} is the typed wrapper for a structural system. It indicates that the wrapped association should be interpreted as a fully assembled equation-system object rather than a generic association.

\subsection*{\texttt{mfgBoundaryData}}

\texttt{mfgBoundaryData[assoc]} wraps boundary-condition data such as entry equations, entry rules, and exit-value rules. It is the typed boundary subrecord stored inside a system.

\subsection*{\texttt{mfgFlowData}}

\texttt{mfgFlowData[assoc]} wraps flow-balance and non-negativity data, including signed-flow expressions and balance equations derived from the auxiliary topology.

\subsection*{\texttt{mfgComplementarityData}}

\texttt{mfgComplementarityData[assoc]} wraps switching-cost consistency results, complementarity alternatives, and optimality inequalities.

\subsection*{\texttt{mfgHamiltonianData}}

\texttt{mfgHamiltonianData[assoc]} wraps the Hamiltonian residual structure, including left-hand and right-hand sides, placeholder cost-current variables, and the general Hamiltonian equation block.

\subsection*{\texttt{mfgSystemQ}, \texttt{mfgBoundaryDataQ}, \texttt{mfgFlowDataQ}, \texttt{mfgComplementarityDataQ}, \texttt{mfgHamiltonianDataQ}}

These predicates recognize the corresponding typed wrappers. Their main role is defensive dispatch. They are used to keep constructors, accessors, and downstream logic attached to the right data shape.

\subsection*{\texttt{makeSystem}}

\texttt{makeSystem[s]} is the main public constructor. In one-argument form it automatically builds the symbolic unknown bundle, then delegates to the two-argument form. In two-argument form, \texttt{makeSystem[s, unk]} builds all four typed subrecords and returns a complete \texttt{mfgSystem}.

This function is the canonical bridge from scenario semantics to symbolic constraints. Typical usage is:

\begin{verbatim}
sys = makeSystem[s];
\end{verbatim}

\subsection*{\texttt{systemData}}

\texttt{systemData[sys]} returns the underlying association. \texttt{systemData[sys,key]} first checks top-level keys and, if necessary, searches inside the typed subrecords before returning \texttt{Missing["KeyAbsent", key]}.

This makes it the standard read interface for every solver and plotting function that consumes system content.

\subsection*{\texttt{systemDataFlatten}}

\texttt{systemDataFlatten[sys]} merges all top-level and typed-subrecord keys into a single flat association. It is primarily a compatibility and inspection helper.

Its low-level importance is that it preserves a bridge to older flat-system access patterns without abandoning the newer typed-record design.

\subsection*{\texttt{getKirchhoffLinearSystem}}

\texttt{getKirchhoffLinearSystem[sys]} extracts the entry-current vector, Kirchhoff matrix, and corresponding variable order. It is a specialized linear-algebra helper for the structural system.

Users working directly on system analysis or certain solver extensions may use this function to access the linear core of the flow-balance structure.

\subsection*{\texttt{getKirchhoffMatrix}}

\texttt{getKirchhoffMatrix[sys]} returns a similar data package but includes a third placeholder slot retained for backward compatibility. It is a compatibility-preserving interface over the same structural idea.

\subsection*{\texttt{consistentSwitchingCosts}}

\texttt{consistentSwitchingCosts[sc][trip]} evaluates whether a particular switching-cost entry obeys a triangle-inequality-style consistency condition relative to alternative intermediate transitions.

This is a diagnostic or analytical helper rather than a standard end-user call. Its main importance is explanatory: it formalizes what the system means by switching-cost consistency.

\subsection*{\texttt{isSwitchingCostConsistent}}

\texttt{isSwitchingCostConsistent[switchingCosts]} checks consistency over the whole switching-cost structure. It accepts both associations and list representations and returns either \texttt{True}, \texttt{False}, or symbolic consistency conditions.

This function directly influences how complementarity alternatives are constructed downstream.

\subsection*{\texttt{altFlowOp}}

\texttt{altFlowOp[j][\{a, b\}]} builds the alternative
\[
j[a, b] = 0 \;\vee\; j[b, a] = 0,
\]
encoding the idea that, on an undirected physical edge, at most one orientation carries flow in the critical complementarity formulation. This is the building block of the disjunctive "alternative zero-flow" block in \texttt{ComplementarityData}.

\subsection*{\texttt{flowSplitting}}

\texttt{flowSplitting[auxTriples][x]} selects triples that begin with the directed pair \texttt{x}. It is used in splitting-balance construction.

\subsection*{\texttt{flowGathering}}

\texttt{flowGathering[auxTriples][x]} selects triples that end with the directed pair \texttt{x}. It is used in gathering-balance construction.

\subsection*{\texttt{switchingCostLookup}}

\texttt{switchingCostLookup[sc]} returns a callable lookup function \texttt{f[r,i,w]} that retrieves the switching cost for a transition, defaulting to zero when absent. This isolates switching-cost representation details from the rest of the system builders.

\subsection*{\texttt{ineqSwitch}}

\texttt{ineqSwitch[u, switching][r, i, w]} builds the local switching optimality inequality at junction $i$ for the transition from $r$ to $w$:
\[
u[w, i] + \mathrm{sc}(r, i, w) - u[r, i] \;\ge\; 0.
\]
It says that diverting through $w$ cannot strictly improve on the value attached to $r$. When the switching cost is supplied as a function rather than a constant, it is applied to the transition flow $j[r, i, w]$ before the inequality is formed, which lets the system accommodate flow-dependent switching penalties.

\subsection*{\texttt{altSwitch}}

\texttt{altSwitch[j, u, switching][r, i, w]} builds the complementarity alternative paired with that inequality:
\[
j[r, i, w] = 0 \;\;\vee\;\; u[w, i] + \mathrm{sc}(r, i, w) - u[r, i] = 0.
\]
Either the transition carries no flow, or the corresponding switching inequality is tight at $i$. The system collects one such alternative per admissible triple $(r, i, w)$; the resulting disjunctive structure is the main source of branching for the symbolic solvers.

\subsection*{\texttt{roundValues}}

\texttt{roundValues[x]} rounds numeric content to a standard precision. The file provides overloads for numbers, rules, lists, and associations. This is a small cleanup utility for making output more stable and readable.

\subsection*{\texttt{exactBoundaryValue} and \texttt{exactBoundaryValues}}

These private helpers convert numeric boundary data into exact values, typically through \texttt{Rationalize[val, 0]}. They reject nonreal or nonnumeric boundary values by returning \texttt{Failure} objects. This exactification step matters because the active symbolic solver pipeline depends on exact structural equations wherever possible.

\subsection*{\texttt{buildBoundaryData}}

\texttt{buildBoundaryData[s, topology]} constructs entry equations, entry rules, exit-value rules, and boundary metadata. It turns raw entry and exit values into exact symbolic constraints and substitutions suitable for the later system and solver pipeline.

\subsection*{\texttt{buildFlowData}}

\texttt{buildFlowData[s, topology, unk]} constructs signed-flow expressions, non-negativity constraints for both \texttt{Js} and \texttt{Jts}, splitting-balance equations, gathering-balance equations, and associated rule forms.

This is one of the most structurally important builders because it transforms topology plus variable families into the conservation laws that define admissible flow behavior.

\subsection*{\texttt{buildComplementarityData}}

\texttt{buildComplementarityData[s, topology, unk]} builds the switching-cost and complementarity layer. It checks switching-cost consistency, constructs alternative zero-flow constraints on undirected edges, builds transition-flow alternatives, assembles switching inequalities by vertex, and forms the optimality complementarity block.

This is the main source of logical disjunction in the final system.

\subsection*{\texttt{boundaryPreservingAutomorphisms}}

\texttt{boundaryPreservingAutomorphisms[scenario]} returns the list of graph automorphisms of the scenario's underlying graph that also preserve the (vertex, value) multisets of \texttt{Entries} and \texttt{Exits}. It is the structural ingredient used to derive symbolic orbit equalities; the function returns the empty list when no nontrivial symmetry survives the boundary data.

\subsection*{\texttt{addSymmetryEqualities}}

\texttt{addSymmetryEqualities[sys, scenario]} returns a new \texttt{mfgSystem} whose \texttt{EqGeneral} block has been extended with the orbit equalities
\[
j[a, b] \;=\; j[\sigma(a), \sigma(b)],
\qquad
u[a, b] \;=\; u[\sigma(a), \sigma(b)],
\]
for every boundary-preserving automorphism $\sigma$ produced by \texttt{boundaryPreservingAutomorphisms}. The augmented system has the same solution set as \texttt{sys}, since any equilibrium is invariant under symmetries of both the graph and the boundary data. The benefit is structural: the solver enters the search with a smaller set of independent unknowns, often by a large factor on symmetric grids and cycles. The function is a no-op when the symmetry group is trivial.

\subsection*{\texttt{addOracleEqualities}}

\texttt{addOracleEqualities[sys, oracleResult]} appends \texttt{j[e] == 0} to \texttt{EqGeneral} for every edge \texttt{e} that the oracle has classified as inactive. It consumes the association returned by \texttt{numericOracleClassify} (in the optional \texttt{numericOracle} subpackage); the function itself is symbolic. Before returning, it runs a \texttt{FindInstance} safety probe on the pruned system: if the over-pruned system has no feasible point, the original \texttt{sys} is returned unchanged. The function is a no-op when the oracle reports no inactive edges or did not converge.

\subsection*{\texttt{buildHamiltonianData}}

\texttt{buildHamiltonianData[s, topology, flowData]} builds the Hamiltonian residual block \texttt{EqGeneral}. For every undirected edge $\{a, b\}$ it records
\[
u[a, b] - u[b, a] + j[a, b] - j[b, a] \;=\; \mathrm{nrhs}_{ab},
\qquad
\mathrm{nrhs}_{ab} \;=\;
\begin{cases}
0 & \alpha_{ab} = 1, \\
m - \operatorname{sign}(m)\, m^{\alpha_{ab}} & \alpha_{ab} \ne 1,
\end{cases}
\]
where $\alpha_{ab}$ is the edge-specific exponent (\texttt{Alpha} globally, overridable per edge by \texttt{EdgeAlpha}). The equation is enforced unconditionally for every edge; in particular, zero-flow edges force $u[a, b] = u[b, a]$, which propagates through the switching inequalities and pins value variables at bypassed exit nodes to their terminal cost. Placeholder cost-current symbols are stored alongside \texttt{EqGeneral} for future non-critical/numeric work but do not participate in the critical solve.

\section{Usage guidance}

End users should normally interact with this module through \texttt{makeSystem} and \texttt{systemData}. More specialized helpers such as \texttt{getKirchhoffMatrix} or switching-consistency functions are most relevant when extending solvers, profiling structural growth, or debugging particular system blocks.

The canonical workflow is:

\begin{verbatim}
unk = makeSymbolicUnknowns[s];
sys = makeSystem[s, unk];
flowBlock = systemData[sys, "FlowData"];
\end{verbatim}

\section*{References}

\begin{enumerate}
\item Source file: \texttt{MFGraphs/systemTools.wl}
\item Upstream dependencies: \texttt{scenarioTools.wl}, \texttt{unknownsTools.wl}
\item Downstream consumer: \texttt{solversTools.wl}
\end{enumerate}

\end{document}
